% Conclusion:
% We did it. This paper rocks and you are lucky to have read it (i.e. brief recap of the entire paper). Also, we’ll do all these other amazing things in the future. 
% To keep going with the analogy, you can think of future work as (potential) academic offspring (credits to James).

\vspace{-0.2cm}
\section{Conclusions} \label{section:conclusions}
\vspace{-0.2cm}
This paper introduces \ours{}, a novel framework for efficient prediction of agentic workflow performance that leverages a multi-view predictive approach. By integrating multi-view graph structures, code semantics, and prompt embeddings into a unified representation, \ours{} captures the diverse characteristics of agentic systems. Moreover, it employs cross-domain unsupervised pretraining to mitigate the challenge of limited labeled data, thereby enhancing generalization across varied tasks. Through comprehensive experiments spanning three domains, \ours{} consistently outperforms strong baselines in predictive accuracy and workflow utility.

\textbf{Limitations and Future Work}. While \ours{} exhibits strong performance, it has certain limitations. The current predictor focuses on binary success metrics, constrained by the available benchmark, which may overlook more nuanced aspects of workflow behavior. Evaluating on new, independently curated agentic benchmarks is an important direction for future work. Additionally, adapting to highly specialized domains may still require some labeled data. Future work includes expanding to multi-objective optimization\,(e.g., balancing accuracy and cost), incorporating richer views such as temporal traces and user feedback, and exploring human-in-the-loop workflows for real-time refinement. These directions aim to make \ours{} more generalizable and interactive in complex, real-world settings.

